% !TEX root =../ECE343_Slides_Rogers.tex


%%%%%%%%%%%%%%%%%%%%%%%%%%%%%%%%%%%%%%%%%%%%%%%%%%%%%%
%%%%%%%%%%%%%%%%%%%%%%%%%%%%%%%%%%%%%%%%%%%%%%%%%%%%%%

%\begin{comment}
\section{Schedule}
\begin{frame}{Schedule}
	\def\myfontsize{footnotesize}
	\setcounter{mylnum}{24}
%\def\myfontsize{normalsize}
\begin{center}
\begin{\myfontsize}
\begin{tabular}{| c | l | c | l |}
\hline
\textbf{Lesson} 		  	& \textbf{Topics}             & \pbox{5cm}{\textbf{Reading}}	& \textbf{Notes}		\\ \hline\hline
\ifnum\lessonNum=\value{mylnum} \rowcolor{green!20!white}\fi
	\themylnum  
	& BVP (Capacitor, Coax), Electric Potential	
	& 5.16
	&  LO 5.3  
\stepcounter{mylnum}\\ \hline \hline
\ifnum\lessonNum=\value{mylnum} \rowcolor{green!20!white}\fi
	\themylnum 	
	& Electric Boundary Conditions
	& 5.18-5.22	
	& LO 5.4			
\stepcounter{mylnum}\\ \hline
\ifnum\lessonNum=\value{mylnum} \rowcolor{green!20!white}\fi
	\themylnum 	
	& BC's/Capacitance			
	& 5.23-5.24
	& 	
\stepcounter{mylnum}\\ \hline
\ifnum\lessonNum=\value{mylnum} \rowcolor{green!20!white}\fi
	\themylnum 	
	& Electrostatic Lab			
	& Lab packet
	& 	
\stepcounter{mylnum}\\ \hline
\rowcolor{gray!20!white}\multicolumn{4}{| c |}{\textbf{Chapter 6: Steady Current \& Conductivity}}	\\ \hline\hline
\ifnum\lessonNum=\value{mylnum} \rowcolor{green!20!white}\fi
	\themylnum 	
	& Current density, Conductivity, Resistance
	& Ch. 6
	&   					
\stepcounter{mylnum}\\ \hline\hline
\rowcolor{gray!20!white}\multicolumn{4}{| c |}{\textbf{Chapter 7: Magnetostatics}}						\\ \hline\hline
\ifnum\lessonNum=\value{mylnum} \rowcolor{green!20!white}\fi
	\themylnum 	
	&  Magnetic Fields, Gauss' Law for Magnetism	
	& 7.1-7.3		
	&LO 7.1-7.2 \textbf{GR 2} 				
\stepcounter{mylnum}\\ \hline
\ifnum\lessonNum=\value{mylnum} \rowcolor{green!20!white}\fi
	\themylnum 	
	& Ampere's Law 	
	& 7.4-7.5, 7.9
	& LO 7.3		
\stepcounter{mylnum}\\ \hline 
\ifnum\lessonNum=\value{mylnum} \rowcolor{green!20!white}\fi
	\themylnum 	
	&  Ampere's Law Applied
	&  7.6-7.8
	& LO 7.3
\stepcounter{mylnum}\\ \hline 
\ifnum\lessonNum=\value{mylnum} \rowcolor{green!20!white}\fi
	\themylnum	
	& Magnetic Boundary Conditions	
	& 7.10-7.14	
	& LO 7.4 		
\stepcounter{mylnum}\\ \hline 
\ifnum\lessonNum=\value{mylnum} \rowcolor{green!20!white}\fi
	\themylnum  
	& Magnetostatics Lab
	& Lab packet
	&  		
\stepcounter{mylnum}\\ \hline \hline
\rowcolor{gray!20!white}\multicolumn{4}{| c |}{\textbf{Chapter 8: Time-Varying Fields}}					\\ \hline\hline
\ifnum\lessonNum=\value{mylnum} \rowcolor{green!20!white}\fi
	\themylnum 	
	& Faraday's Law, Displacement Current, Maxwell Eqns
	& 8.3, 8.8, 8.9
	& LO 8.1
\stepcounter{mylnum}\\ \hline 
\ifnum\lessonNum=\value{mylnum} \rowcolor{green!20!white}\fi
	\themylnum	
	& Faraday's Law, Displacement Current, Maxwell Eqns 
	&  8.3, 8.8, 8.9
	& LO 8.2-8.3
\stepcounter{mylnum}\\ \hline 

\rowcolor{gray!20!white}\multicolumn{4}{| c |}{\textbf{Chapter 9: Plane Waves in Lossless Media}}					\\ \hline\hline
\ifnum\lessonNum=\value{mylnum} \rowcolor{green!20!white}\fi
	\themylnum	
	& Maxwell's Equations in lossless, source-free region		
	& 9.1-9.2
	& LO 9.1	 
\stepcounter{mylnum}\\ \hline
\ifnum\lessonNum=\value{mylnum} \rowcolor{green!20!white}\fi
	\themylnum	
	& Maxwell's Equations in lossless, source-free region		
	& 9.1-9.2
	& LO 9.1	 
\stepcounter{mylnum}\\ \hline
%\ifnum\lessonNum=\value{mylnum} \rowcolor{green!20!white}\fi
%	\themylnum	
%	& Flex day
%	& 
%	& \textbf{GR 3}			 					
%\stepcounter{mylnum}\\ \hline 
\ifnum\lessonNum=\value{mylnum} \rowcolor{green!20!white}\fi
	\themylnum	
	& Uniform Plane Waves \& Polarization
	& 9.4-9.6
	& LO 9.2-9.3			
\stepcounter{mylnum}\\ \hline 
\ifnum\lessonNum=\value{mylnum} \rowcolor{green!20!white}\fi
	\themylnum	
	& Normal Incidence
	& Vol 2, 5.6-5.7
	&  LO 9.4-9.5
\stepcounter{mylnum}\\ \hline 
\ifnum\lessonNum=\value{mylnum} \rowcolor{green!20!white}\fi
	\themylnum	
	& Oblique Incidence			
	& Vol 2, 5.6-5.7
	& LO 9.5
\stepcounter{mylnum}\\ \hline 
\ifnum\lessonNum=\value{mylnum} \rowcolor{green!20!white}\fi
	\themylnum	
	& Oblique Incidence			
	& Vol 2, 5.6-5.7
	& LO 9.5
\stepcounter{mylnum}\\ \hline 
\ifnum\lessonNum=\value{mylnum} \rowcolor{green!20!white}\fi
	\themylnum	
	& Incidence Lab			
	& TBD
	& LO 9.5
\stepcounter{mylnum}\\ \hline 
%\ifnum\lessonNum=\value{mylnum} \rowcolor{green!20!white}\fi
%	\themylnum	
%	& Wrapup
%	& 
%	& 		 					
%\stepcounter{mylnum}\\ \hline 

\end{tabular}			
\end{\myfontsize}
\end{center}

\end{frame}

%%%%%%%%%%%%%%%%%%%%%%%%%%%%%%%%%%%%%%%%%%%%%%%%%%%%%
%%%%%%%%%%%%%%%%%%%%%%%%%%%%%%%%%%%%%%%%%%%%%%%%%%%%%

\def\obja{ I understand the process of decoupling and finding solutions for Maxwell's equations for dynamic fields.}
\def\objb{ I can articulate what a Uniform Plane Wave (UPW) is and describe key properties that relate the electric and magnetic field. }
\def\objc{ I can calculate the electric and magnetic fields for a UPW.}
\def\objd{ I can articulate the concept of polarization and determine the polarization for various UPWs.}
\def\obje{ I can determine the reflection and transmission of UPW at a boundary for Transverse Electric (TE) and Transverse Magnetic (TM) waves, using boundary conditions.}
\begin{frame}{Ch 9 Learning Objectives}
	\begin{textblock}{14.5}(0.5,2.5)
	\begin{myblock}[0.95]{}
	\begin{small}
	\begin{enumerate}[{9}.1]
		\ifnum\objNum=1
			\item \alert{ \obja }
		\else 
			\item \obja
		\fi 
		\ifnum\objNum=23
			\item \alert{ \objb }
			\item \alert{ \objc }
		\else 
			\item \objb
			\item \objc
		\fi 
		\ifnum\objNum=4
			\item \alert{ \objd }
		\else 
			\item \objd
		\fi 
		\ifnum\objNum=5
			\item \alert{ \obje }
		\else 
			\item \obje
		\fi 
	\end{enumerate}
	\end{small}
	\end{myblock}
	\end{textblock}
\end{frame}

%%%%%%%%%%%%%%%%%%%%%%%%%%%%%%%%%%%%%%%%%%%%%%%%%%%%%%
%%%%%%%%%%%%%%%%%%%%%%%%%%%%%%%%%%%%%%%%%%%%%%%%%%%%%%
\begin{frame}{Types of Waves: By Shape of Wave Front}
	\hspace{-1.25cm}\begin{tikzpicture}
	\definecolor{myblue}{RGB}{187,224,227}
	\draw[fill=blue!10!white, draw=white] (0,0) rectangle ++(6.65,10);
	\end{tikzpicture}
	
	\begin{textblock}{7.5}(0.5,3.5) % left blue bar space
	Surfaces of constant phase can be:
	\begin{itemize}
		\setlength\itemsep{14pt}
		\item Spherical
			
			Appropriate when \handoutgap{0.5in}{source} is  \handoutgap{0.5in}{small} relative to  \handoutgap{0.5in}{domain size}
			
		\item Cylindrical
		
		Appropriate when \handoutgap{0.5in}{source} is shaped like a \handoutgap{0.5in}{line}.
		
		\item Planar
		
		Appropriate when \handoutgap{0.5in}{wave} is observed \handoutgap{0.5in}{locally} or when wave is intentionally guided to be \handoutgap{0.5in}{planar}. 
	\end{itemize}
	\end{textblock}
	\begin{textblock}{8}(8.5,3)
	
		\hspace{0.5in}\includegraphics[scale=0.15]{./ECE343/spherical_9-1}
		
		\hspace{0.5in}\includegraphics[scale=0.15]{./ECE343/cylindrical_9-2}
		
		\includegraphics[scale=0.15]{./ECE343/planar_9-3} 		
	\end{textblock}

	\begin{textblock}{5}(12,12.5)
		\includegraphics[scale=0.15]{./ECE343/parabolic_reflector}
	\end{textblock}
\end{frame}



%%%%%%%%%%%%%%%%%%%%%%%%%%%%%%%%%%%%%%%%%%%%%
%%%%%%%%%%%%%%%%%%%%%%%%%%%%%%%%%%%%%%%%%%%%%
%\begin{frame}{Uniform Plane Waves}
%\begin{textblock}{7.5}(0.5,3.5) 
%% Point source
%\begin{tikzpicture}
%  \fill[myblue] (0,0) circle (\p);
%  \foreach \i in {1,...,4}{
%    \draw[myblue,thick] (0,0) circle (\r*\i);
%  }
%  \foreach \a in {30,90,...,330}{
%    \draw[myarr,myblue!80!black] (\a:\r*3.6) --++ (\a:1.2*\r);
%  }
%\end{tikzpicture}
%\end{textblock}
%\end{frame}

%%%%%%%%%%%%%%%%%%%%%%%%%%%%%%%%%%%%%%%%%%%%%%
%%%%%%%%%%%%%%%%%%%%%%%%%%%%%%%%%%%%%%%%%%%%%%
%
\begin{frame}{Uniform Plane Waves}
	\begin{textblock}{7.5}(0.5,3.5) 
	Assumptions:
	\begin{itemize}
		\item Simple Media:
		\begin{itemize}
			\setlength\itemsep{8pt}
			\item \handoutgap{1in}{Linear}
			\item \handoutgap{1in}{Isotropic}
			\item \handoutgap{1in}{Homogeneous}
			\item \handoutgap{1in}{Non-dispersive}
		\end{itemize}
		\item Magnitude and phase are constant in a plane perpendicular to direction of propagation (assume $z$)
		\visible<0|handout:1>{
		\alert{
		\eq{
			\partd{\phasor{E}}{x} = 	\partd{\phasor{E}}{y} = \partd{\phasor{H}}{x} = \partd{\phasor{H}}{y} = 0
		}  % end eq
		} % end alert
		} % end visible
	\end{itemize}
	\end{textblock}

\begin{textblock}{7.5}(8.25,2.5)
Which means:

\visible<0|handout:1>{
\alert{
\eq{
	\partd{}{x}\phasor{E}_x = \partd{}{x}\phasor{E}_y = \partd{}{x}\phasor{E}_z = 0\\
	\partd{}{x}\phasor{H}_x = \partd{}{x}\phasor{H}_y = \partd{}{x}\phasor{H}_z = 0
} % end eq
} % end alert
} % end visible

Now, the wave equation can be written as: 
\begin{small}
\visible<0|handout:1>{
\alert{
\eq{
	&\lb \partd{^2}{x^2} + \partd{^2}{y^2} + \partd{^2}{z^2} \rb 
	\lp \xhat \phasor{E}_x + \yhat \phasor{E}_y  + \zhat \phasor{E}_z \rp  \\
	&+ \gamma^2 \lp \xhat \phasor{E}_x + \yhat \phasor{E}_y  + \zhat \phasor{E}_z \rp = 0 
} % end eq
} % end alert
} % end visible
\end{small}
\end{textblock}
\end{frame}

%%%%%%%%%%%%%%%%%%%%%%%%%%%%%%%%%%%%%%%%%%%%%%%%%%%%%%%%%%%%%%%%%%
%%%%%%%%%%%%%%%%%%%%%%%%%%%%%%%%%%%%%%%%%%%%%%%%%%%%%%%%%%%%%%%%%%
\begin{frame}{Uniform Plane Waves: Fields}
\begin{textblock}{7.5}(0.5,3) 
Each component separately must equal zero, allowing us to write 3 equations:
\begin{small}
\visible<0|handout:1>{
\alert{
\eq{
	\lb \partd{^2}{x^2} + \partd{^2}{y^2} + \partd{^2}{z^2} \rb \phasor{E}_{x,y,z} + \gamma^2 \phasor{E}_{x,y,z} = 0 
} % end eq
} % end alert
} % end visible
\end{small}
So, the UPW conditions reduce these equations to:
\vspace{-12pt}
\begin{myblock}[0.9]{}
\begin{small}
\visible<0|handout:1>{
\alert{
\vspace{-12pt}
\eq{
	\partd{^2}{z^2}\phasor{E}_{x,y,z} + \gamma^2 \phasor{E}_{x,y,z} \\
	} % end eq
	\vspace{-12pt}
	} % end alert
	} % end visible
\end{small}
\end{myblock}
\end{textblock}


\begin{textblock}{7.5}(8.5,3)
We can make one more simplification, by considering Ampere's Law:
%\vspace{-8pt}
\begin{small}
\visible<0|handout:1>{
\alert{
\eq{
	\curl{\phasor{H}} &= \lp \sigma +  j\omega \varepsilon\rp\phasor{E}\\
	\zhat \mycdot \lp \curl{\phasor{H}} \rp &= \zhat\mycdot\lb \lp \sigma +  j\omega \varepsilon\rp\phasor{E} \rb \\
	\partd{\phasor{H}_x}{y} - \partd{\phasor{H}_y}{x}  &= \lp \sigma +  j\omega \varepsilon\rp \phasor{E}_{z} \\
	0-0 &= \lp \sigma +  j\omega \varepsilon\rp\phasor{E}_{z} \\
	0 &= \phasor{E}_{z}
} % end eq
} % end alert
} % end visible
\end{small}
%\only<1|handout:0>{\vspace{-32pt}}
\vspace{-24pt}
\begin{myblock}[0.9]{}
\begin{small}
\vspace{-12pt}
\visible<0|handout:1>{
\alert{
\eq{
	\handoutgap{0pt}{\partd{^2}{z^2}\phasor{E}_{x,y} + \gamma^2 \phasor{E}_{x,y}} 
}	% end eq
\vspace{-12pt}
} % end alert
} % end visible
\end{small}
\end{myblock}
\end{textblock}
\end{frame}

%%%%%%%%%%%%%%%%%%%%%%%%%%%%%%%%%%%%%%%%%%%%%%%%%%%%%%%%%%%%%%%%%%
%%%%%%%%%%%%%%%%%%%%%%%%%%%%%%%%%%%%%%%%%%%%%%%%%%%%%%%%%%%%%%%%%%
\begin{frame}{Uniform Plane Waves: Fields}
\begin{textblock}{7.5}(0.5,3.5) 
We know how to find solutions for the 2nd order ODEs (right?!)
\begin{small}
\visible<0|handout:1>{
\vspace{-12pt}
\alert{
\eq{
	&\partd{^2}{z^2}\phasor{E}_{x} + \gamma^2 \phasor{E}_{x}  = 0\\
	&\;\implies\;
	\phasor{E}_x (z) = E^{+}_{x0} e^{-j\gamma z}  + E^{-}_{x0} e^{j\gamma z} \\
	&\partd{^2}{z^2}\phasor{E}_{y} + \gamma^2 \phasor{E}_{y} = 0 \\
	&\;\implies\;
	\phasor{E}_y (z) = E^{+}_{y0} e^{-j\gamma z}  + E^{-}_{y0} e^{j\gamma z} \\
	&\phasor{E} = \xhat \phasor{E}_x + \yhat \phasor{E}_y \\
	&\phasor{H}_x (z) = H^{+}_{x0} e^{-j\gamma z}  + H^{-}_{x0} e^{j\gamma z} \\
	&\phasor{H}_y (z) = H^{+}_{y0} e^{-j\gamma z}  + H^{-}_{y0} e^{j\gamma z} \\
	&\phasor{H} = \xhat \phasor{H}_x + \yhat \phasor{H}_y\\
	&\gamma = \alpha + j\beta = \sqrt{j\omega\mu\lp\sigma + j\omega\varepsilon\rp } 
} % end eq
} % end alert
} % end visible
\end{small}

\end{textblock}


\begin{textblock}{7.5}(8.5,3.5)
Now, back to time domain...multiply by \handoutgap{1in}{$e^{j\omega t}$} and \handoutgap{1.5in}{take the real part}.

\begin{small}
\visible<0|handout:1>{
\alert{
\eq{
	E_{x,y} (z,t) &= \text{Re}\left\{ E^{+}_{x0, y0} e^{-j\beta z}e^{j\omega t} \right.\\
	&\; \left. + E^{-}_{x0,y0} e^{j\beta z}e^{j\omega t}\right\} \\
	E_{x,y} (z,t) &= E^{+}_{x0,y0} e^{-\alpha z}\cosp{\omega t - \beta z} \\
	&\; + E^{-}_{x0,y0}e^{\alpha z} \cosp{\omega t + \beta z}
} % end eq
} % end alert
} % end visible
\end{small}

\end{textblock}
\end{frame}

%%%%%%%%%%%%%%%%%%%%%%%%%%%%%%%%%%%%%%%%%%%%%%%%%%%%%%%%%%%%%%%%%%
%%%%%%%%%%%%%%%%%%%%%%%%%%%%%%%%%%%%%%%%%%%%%%%%%%%%%%%%%%%%%%%%%%
%\begin{frame}{Uniform Plane Waves: Fields Summary}
%\begin{textblock}{7.75}(0.5,2.5) 
%\begin{myblock}[0.95]{General}
%\begin{small}
%\eq{
%	&\vec{E} (z,t) = \xhat E_{x} + \yhat E_{y} (z,t)\\
%	&E_{x,y} (z,t) = E^{+}_{x0,yo} e^{-\alpha z}\cosp{\omega t - \beta z} \\
%	&\qquad\qquad + E^{-}_{x0,y0} e^{\alpha z}\cosp{\omega t + \beta z}  \\[12pt]
%	&\vec{H} (z,t) = \xhat H_{x} + \yhat H_{y} (z,t)\\
%	&H_{x,y} (z,t) = H^{+}_{x0,yo} e^{-\alpha z}\cosp{\omega t - \beta z} \\
%	&\qquad\qquad + H^{-}_{x0,y0} e^{\alpha z}\cosp{\omega t + \beta z}  \\
%}
%\end{small}
%\end{myblock}
%\end{textblock}
%
%\begin{textblock}{7.25}(8.25,2.5) 
%\begin{myblock}[0.95]{Lossless}
%\begin{small}
%\visible<0|handout:1>{
%\eq{
%	&\vec{E} (z,t) = \xhat E_{x} + \yhat E_{y} (z,t)\\
%	&E_{x,y} (z,t) = E^{+}_{x0,y0} \cosp{\omega t - \beta z} \\
%	&\qquad\qquad + E^{-}_{x0,y0} \cosp{\omega t + \beta z}  \\[12pt]
%	&\vec{H} (z,t) = \xhat H_{x} + \yhat H_{y} (z,t)\\
%	& H_{x,y} (z,t) = H^{+}_{x0,y0}\cosp{\omega t - \beta z} \\
%	&\qquad\qquad + H^{-}_{x0,y0} \cosp{\omega t + \beta z}  \\[12pt]
%} % end eq
%} % end visible
%\end{small}
%\end{myblock}
%\end{textblock}
%\end{frame}
%%%%%%%%%%%%%%%%%%%%%%%%%%%%%%%%%%%%%%%%%%%%%%%%%%%%%%%%%%%%%%%%%%
%%%%%%%%%%%%%%%%%%%%%%%%%%%%%%%%%%%%%%%%%%%%%%%%%%%%%%%%%%%%%%%%%%

\begin{frame}{UPW Summary}
\begin{textblock}{8}(0.5,2)
\begin{myblock}[0.95]{Maxwell's Equations}
\vspace{-12pt}
\eq{
&\text{Gauss}\longrightarrow\div{\phasor{E}} = \dfrac{\rho_v}{\varepsilon} \\
&\text{Gauss}\longrightarrow\div{\phasor{H}} = 0 \\
&\text{Ampere}\longrightarrow\curl{\phasor{H}} = \lp \sigma + j\omega\varepsilon \rp \phasor{E} \\
&\text{Faraday}\longrightarrow\curl{\phasor{E}} = -j\omega\mu\phasor{H} 
}
\vspace{-12pt}
\end{myblock}
\end{textblock}
\begin{textblock}{6}(9,2)
\begin{myblock}[0.9]{Helmholtz}
\vspace{-12pt}
\eq{
&\nabla^2 \phasor{E} - \gamma^2 \phasor{E} = 0 \\
&\nabla^2 \phasor{H} - \gamma^2 \phasor{H} = 0 \\
&\gamma^2 = j\omega\mu\lp\sigma + j\omega\varepsilon\rp \\[-4pt]
}
\vspace{-12pt}
\end{myblock}
\end{textblock}
\begin{textblock}{14.5}(0.5,9.75)
\begin{myblock}[0.99]{Field Solutions}
\vspace{-18pt}
\eq{
&E_{x,y} (z,t) = E^{+}_{x0,y0} e^{-\alpha z}\cosp{\omega t - \beta z} 
+ E^{-}_{x0,y0} e^{\alpha z}\cosp{\omega t + \beta z}  \\
&H_{x,y} (z,t) = H^{+}_{x0,y0} e^{-\alpha z}\cosp{\omega t - \beta z} 
+ H^{-}_{x0,y0} e^{\alpha z}\cosp{\omega t + \beta z}  
}
\vspace{-18pt}
\end{myblock}
\end{textblock}
\end{frame}

%%%%%%%%%%%%%%%%%%%%%%%%%%%%%%%%%%%%%%%%%%%%%%%%%%%%%%%%%%%%%%%%%%
%%%%%%%%%%%%%%%%%%%%%%%%%%%%%%%%%%%%%%%%%%%%%%%%%%%%%%%%%%%%%%%%%%

\begin{frame}{Uniform Plane Waves: Interpretation}
\begin{textblock}{14.5}(0.5,2.25)

\begin{myblock}[0.98]{Observations}
	\begin{itemize}
		\item For a UPW, the electric and magnetic field \handoutgap{1in}{vectors} are \handoutgap{1in}{perpendicular} to propagation.
%		\item For a UPW, the wave may \handoutgap{1in}{propagate} only \handoutgap{1in}{perpendicular} to the plane
		\item The direction of propagation is \handoutgap{1in}{perpendicular} to the \handoutgap{1in}{electric field} and \handoutgap{1in}{magnetic field}
		\item The \handoutgap{1in}{direction} of the \handoutgap{1in}{electric field} and \handoutgap{1in}{magnetic field} are related
	\end{itemize}

\end{myblock}


\end{textblock}

\begin{textblock}{7.5}(0.5,11.5)
\scalebox{0.5}{
% Electromagnetic wave - colored
\begin{tikzpicture}[x=(-15:0.9), y=(90:0.9), z=(-150:1.1),
                    line cap=round, line join=round,
                    axis/.style={black, thick,->},
                    vector/.style={>=stealth,->}]
  \large
  \def\A{1.5}
  \def\nNodes{5} % use even number
  \def\nVectorsPerNode{8}
  \def\N{\nNodes*40}
  \def\xmax{\nNodes*pi/2*1.01}
  \pgfmathsetmacro\nVectors{(\nVectorsPerNode+1)*\nNodes}
  \def\vE{{\color{Ecol}\mathbf{E}}}
  \def\vB{{\color{Bcol}\mathbf{B}}}
  
  \def\drawENode{ % draw E node and vectors with some offset
    \draw[Ecol,very thick,variable=\t,domain=\iOffset*pi/2:(\iOffset+1)*pi/2*1.01,samples=40]
      plot (\t,{\A*sin(\t*360/pi)},0);
    \foreach \k [evaluate={\t=\k*pi/2/(\nVectorsPerNode+1);
                           \angle=\k*90/(\nVectorsPerNode+1);}]
                in {1,...,\nVectorsPerNode}{
      \draw[vector,EVcol]  (\iOffset*pi/2+\t,0,0) -- ++(0,{\A*sin(2*\angle+\iOffset*180)},0);
    }
  }
  \def\drawBNode{ % draw B node and vectors with some offset
    \draw[Bcol,very thick,variable=\t,domain=\iOffset*pi/2:(\iOffset+1)*pi/2*1.01,samples=40]
      plot (\t,0,{\A*sin(\t*360/pi)});
    \foreach \k [evaluate={\t=\k*pi/2/(\nVectorsPerNode+1);
                           \angle=\k*90/(\nVectorsPerNode+1);}]
                in {1,...,\nVectorsPerNode}{
      \draw[vector,Bcol!50]  (\iOffset*pi/2+\t,0,0) -- ++(0,0,{\A*sin(2*\angle+\iOffset*180)});
    }
  }
  
  % MAIN AXES
  \draw[axis] (0,0,0) -- ++(\xmax*1.1,0,0) node[right] {$z$};
  \draw[axis] (0,-\A*1.2,0) -- (0,\A*1.2,0) node[right] {$x$};
  \draw[axis] (0,0,-\A*1.2) -- (0,0,\A*1.2) node[left] {$y$};
  
  % SMALL AXES
  \def\xOffset{{(\nNodes-2)*pi/2}}
  \def\yOffset{\A*1.1}
  \def\zOffset{\A*1.1}
  \draw[axis,very thick,vcol] (\xOffset,\yOffset,-\zOffset) -- ++(\A*0.6,0,0) node[right,align=center] {$\vb{v}$}; %\\propagation
  \draw[axis,very thick,Ecol]  (\xOffset,\yOffset,-\zOffset) -- ++(0,\A*0.6,0) node[right] {$\vb{E}$};
  \draw[axis,very thick,Bcol]   (\xOffset,\yOffset,-\zOffset) -- ++(0,0,\A*0.6) node[left] {$\vb{B}$};
  
  % draw (anti-)nodes
  \foreach \iNode [evaluate={\iOffset=\iNode-1;}] in {1,...,\nNodes}{
    \ifodd\iNode \drawBNode \drawENode % E overlaps B
    \else        \drawENode \drawBNode % B overlaps E
    \fi
  }

\end{tikzpicture}
} % end scalebox
\end{textblock}

\begin{textblock}{7.5}(8,11.5)
\scalebox{0.75}{
% Electromagnetic wave SHORT - colored
\begin{tikzpicture}[x=(-15:0.9), y=(90:0.9), z=(-150:1.1),
                    line cap=round, line join=round,
                    axis/.style={black, thick,->},
                    vector/.style={>=stealth,->}]
  \large
  \def\A{1.5}
  \def\ymax{1.8}
  \def\nNodes{1} % use even number
  \def\nVectorsPerNode{8}
  \def\N{\nNodes*40}
  \def\xmax{\nNodes*pi/2*1.01}
  \pgfmathsetmacro\nVectors{(\nVectorsPerNode+1)*\nNodes}
  
  % MAIN AXES
  \draw[axis] (0,0,0) -- ++(\xmax*1.4,0,0) node[right] {$z$};
  \draw[axis] (0,-\ymax*0.6,0) -- (0,\ymax,0) node[right] {$x$};
  \draw[axis] (0,0,-\ymax*0.7) -- (0,0,\ymax) node[left] {$y$};
  
  \draw[Ecol,very thick,variable=\t,domain=0:pi/2*1.01,samples=40]
    plot (\t,{\A*sin(\t*360/pi)},0);
  \foreach \k [evaluate={\t=\k*pi/2/(\nVectorsPerNode+1);
                         \angle=\k*90/(\nVectorsPerNode+1);}]
              in {1,...,\nVectorsPerNode}{
    \draw[vector,EVcol]  (\t,0,0) -- ++(0,{\A*sin(2*\angle)},0);
  }
  \draw[Bcol,very thick,variable=\t,domain=0:pi/2*1.01,samples=40]
    plot (\t,0,{\A*sin(\t*360/pi)});
  \foreach \k [evaluate={\t=\k*pi/2/(\nVectorsPerNode+1);
                         \angle=\k*90/(\nVectorsPerNode+1);}]
              in {1,...,\nVectorsPerNode}{
    \draw[vector,Bcol!50]  (\t,0,0) -- ++(0,0,{\A*sin(2*\angle)});
  }
  
  % VECTOR
  \draw[->,very thick,vcol] (\A*1.1,\A*0.8,0) --++ (\A*0.5,0,0) node[right] {$\vb{v}$};
  
\end{tikzpicture}
} % End scalebox
\end{textblock}
\end{frame}


%%%%%%%%%%%%%%%%%%%%%%%%%%%%%%%%%%%%%%%%%%%%%
%%%%%%%%%%%%%%%%%%%%%%%%%%%%%%%%%%%%%%%%%%%%%
\section{Interpretation}
\begin{frame}{Chapter 3: Lesson 5}
	\begin{textblock}{8}(2.5,3.5)
		\vspace{-12pt}
			\eq{
				V(z,t) &=
				\underbrace{\alert<2>{V^{+}_0} \alert<3>{e^{-\alpha z}}\alert<4>{\cosp{\omega t - \beta z}}}_{\color{magenta}{\text{Forward Wave}}}
				+   \underbrace{\alert<2>{V^{-}_0} \alert<3>{e^{\alpha z}}\alert<4>{\cosp{\omega t + \beta z}}}_{\text{\color{green!50!black}{Reverse Wave}}}
				} % end eq 
			
	\end{textblock}

	\begin{textblock}{9}(0.5,6.5)
		\begin{itemize}
			\setlength\itemsep{12pt}
			\item<2-| alert@2> $V^{\pm}_{0} \longrightarrow$ Amplitude coefficients
			\item<3-| alert@3> $e^{\pm \alpha z} \longrightarrow$ Attenuation term
			\item<4-| alert@4> $\cosp{\omega t \pm \beta z} \longrightarrow$Time and space-varying wave
		\end{itemize}
	\end{textblock}
	
	\begin{textblock}{9}(8.5,6.5)
	\only<5>{
	\scalebox{0.75}{
	\begin{tikzpicture}
		\begin{axis}[domain=0:60e-3, 
			samples=200, 
			xmin=0, 
			xmax=65e-3, 
			ymax=4,
			ymin=-4,
	%		grid=major, 
			title={Forward Wave},
			scaled x ticks=base 10:3,
			xtick scale label code/.code={},
			minor x tick num = 1,
			minor y tick num = 1,
	%		ytick={-400,-300,-200,-100,0,100,200,300,400},
	%		xtick={-15e-3,-10e-3,-5e-3,0,5e-3,10e-3,15e-3,20e-3},
			yticklabel={\empty},
			xticklabel={\empty},
			xlabel=$z$,
			ylabel={$V(z,t)$},
	%		legend entries={$f=60$ Hz, $f=120$ Hz},
	%		legend pos=outer north east,
			axis y line = middle,
			axis x line = middle,
			x label style={xshift=0.5cm, yshift=-0.5cm},
			y label style={xshift=-0.5cm, yshift=0.75cm},
			clip=false,
			]
			\addplot [mark=none, orange!50!black, dashed] {cos(360*50*x)} 
				node[rounded corners,draw=orange!50!black,dashed,text=black,rectangle callout, fill=none, callout absolute pointer={(axis cs:12e-3,-1.1)}] at (axis cs: 20e-3, -2) {$\cosp{\omega t - \beta z}$};
			\addplot [dashed, magenta] {2*exp(-50*x)} 
				node[rounded corners,draw=magenta,dashed,text=black,rectangle callout, fill=none, callout absolute pointer={(axis cs:28e-3,0.5)}] at (axis cs: 25e-3, 2) {$e^{-\alpha z}$};
			\addplot[blue] {2*exp(-50*x)*cos(360*50*x)} 
				node[rounded corners,draw=blue,solid,text=black,rectangle callout, fill=none, callout absolute pointer={(axis cs:50e-3,-0.1)}] at (axis cs: 50e-3, 1.75) {$e^{-\alpha z}\cosp{\omega t - \beta z}$};

		\end{axis}
	\end{tikzpicture}
	} % end scalebox
	} % end only
	
	\only<6>{
	\scalebox{0.75}{
	\begin{tikzpicture}
		\begin{axis}[domain=0:60e-3, 
			samples=200, 
			xmin=0, 
			xmax=65e-3, 
			ymax=10,
			ymin=-10,
	%		grid=major, 
			title={Reverse Wave},
			scaled x ticks=base 10:3,
			xtick scale label code/.code={},
			minor x tick num = 1,
			minor y tick num = 1,
	%		ytick={-400,-300,-200,-100,0,100,200,300,400},
	%		xtick={-15e-3,-10e-3,-5e-3,0,5e-3,10e-3,15e-3,20e-3},
			yticklabel={\empty},
			xticklabel={\empty},
			xlabel=$z$,
			ylabel={$V(z,t)$},
	%		legend entries={$f=60$ Hz, $f=120$ Hz},
	%		legend pos=outer north east,
			axis y line = middle,
			axis x line = middle,
			x label style={xshift=0.5cm, yshift=-0.5cm},
			y label style={xshift=-0.5cm, yshift=0.75cm},
			clip=false,
			]
			\addplot [mark=none, yellow!50!black, dashed] {cos(360*50*x)};
				%node[rounded corners,draw=blue,solid,text=black,rectangle callout, fill=none, callout absolute pointer={(axis cs:12e-3,-1.1)}] at (axis cs: 20e-3, -2) {$\cosp{\omega t + \beta z}$};
			\addplot [dashed, magenta] {0.5*exp(50*x)} ;
%				node[rounded corners,draw=magenta,solid,text=black,rectangle callout, fill=none, callout absolute pointer={(axis cs:20e-3,1)}] at (axis cs: 25e-3, 2) {$e^{\alpha z}$};
			\addplot[blue] {0.5*exp(50*x)*cos(360*50*x)} ;
				%node[rounded corners,draw=orange!50!black,solid,text=black,rectangle callout, fill=none, callout absolute pointer={(axis cs:50e-3,-0.1)}] at (axis cs: 50e-3, 1.75) {$e^{\alpha z}\cosp{\omega t + \beta z}$};

		\end{axis}
	\end{tikzpicture}
	} % end scalebox
	} % end only
	
	\end{textblock}
	
	% do a graph
\end{frame}

%%%%%%%%%%%%%%%%%%%%%%%%%%%%%%%%%%%%%%%%%%%%%%%%%%%%%%%%%%%%%%%%%
%%%%%%%%%%%%%%%%%%%%%%%%%%%%%%%%%%%%%%%%%%%%%%%%%%%%%%%%%%%%%%%%%
\section{Characteristics}
\begin{frame}{Uniform Plane Waves: Characteristics}
\begin{textblock}{7.5}(0.5,3) 
Assume an $\xhat\;$ Electric field... 
\begin{small}
\visible<0|handout:1>{
\alert{
\eq{
	&\phasor{E} (z) = \xhat \phasor{E}_x (z) \\
	&\curl{\phasor{E}} = -j\omega\mu\phasor{H} \\
	&\xhat\lp 0 \rp  - \yhat \lp 0 - \partd{\phasor{E}_x}{z} \rp + \zhat \lp 0 - \partd{\phasor{E}_x}{y} \rp \\[8pt]%= -j\omega\mu\phasor{H} \\ 
	& \implies \yhat\partd{\phasor{E}_x}{z}
	=  -j\omega\mu\phasor{H}  \\
	&\partd{\phasor{E}_x}{z} = \partd{}{z} \lp E^{+}_{x0} e^{-\gamma z} \rp + \partd{}{z} \lp E^{-}_{x0} e^{\gamma z} \rp \\
	&\implies \yhat\lb -\gamma E^{+}_{x0} e^{-\gamma z} + \gamma E^{-}_{x0} e^{\gamma z} \rb =  -j\omega\mu\phasor{H}  \\[8pt]
	&\implies \phasor{H} = \yhat\lb \dfrac{ E^{+}_{x0}}{j\omega\mu} e^{-\gamma z} -  \dfrac{ E^{-}_{x0}}{j\omega\mu} e^{\gamma z}\rb \\[8pt]
	&\implies \phasor{H} = \yhat \lb  H^{+}_{y0}e^{-\gamma z} + H^{-}_{y0}e^{\gamma z} \rb\\
} % end eq
} % end alert
} % end visible
\end{small}
\end{textblock}


\begin{textblock}{0.5}(8.75,3)
	{\color{Blue4}{\rule{1pt}{2.8in}}}
\end{textblock}


\begin{textblock}{6.5}(9,3)
Intrinsic impedance is like characteristic impedance...
\begin{small}
\visible<0|handout:1>{
\vspace{-8pt}
\alert{
	\eq{
		\eta &= \dfrac{E^{+}_{x0}}{H^{+}_{y0}} 
			= \dfrac{E^{+}_{x0}}{\frac{E^{+}_{x0}}{j\omega\mu}} 
			= \dfrac{j\omega\mu}{\gamma}  \\
		 &= \dfrac{j\omega\mu}{\sqrt{j\omega\mu\lp\sigma + j\omega\varepsilon\rp}} 
		  = \sqrt{\dfrac{j\omega\mu}{\sigma + j\omega\varepsilon}}
	} % end eq
	} % end alert
	} % end visible	
\end{small}
$\phasor{E}$, $\phasor{H}$, and propagation ($\hat{a}_p$) are \textbf{generally} related:
\begin{small}
\visible<0|handout:1>{
\vspace{-8pt}
\alert{
\eq{
	&\phasor{E}  = -\eta \hat{a}_{p} \times \phasor{H} \\
	&\phasor{H} = \dfrac{1}{\eta} \hat{a}_{p} \times \phasor{E}
} % end eq
} % end alert
} % end visible
\end{small}
\end{textblock}
\end{frame}
%\begin{myblock}[0.9]{Electric Field}
%
%
%\end{myblock}


%%%%%%%%%%%%%%%%%%%%%%%%%%%%%%%%%%%%%%%%%%%%%%%%%%%%%%%%%%%%%%%%%%
%%%%%%%%%%%%%%%%%%%%%%%%%%%%%%%%%%%%%%%%%%%%%%%%%%%%%%%%%%%%%%%%%%
\section{General UPW}
\begin{frame}{Uniform Plane Waves: General}
\begin{textblock}{14.5}(0.5,3) 
\begin{myblock}[0.9]{Summary}
\vspace{-8pt}
\eq{
	&\vec{E}(z,t)  =  \xhat {E}_x (z,t) + \yhat {E}_y (z,t) \\[4pt]
	&E_{x,y} (z,t) = E^{\pm}_{x0,y0} e^{\mp \alpha z}\cosp{\omega t \mp \beta z} \\[4pt]
	&\phasor{H}(z,t)  =  \xhat {H}_x (z,t) + \yhat {H}_y (z,t) \\[4pt]
	&H_{x,y} (z,t) = H^{\pm}_{x0,y0} e^{\mp\alpha z}\cosp{\omega t \mp \beta z} \\[4pt]
	&\phasor{E}  = -\eta \hat{a}_{p} \times \phasor{H} \qquad
	\phasor{H} = \dfrac{1}{\eta} \hat{a}_{p} \times \phasor{E}
	&\eta = \sqrt{\dfrac{j\omega\mu}{\sigma + j\omega\varepsilon}}
}

\end{myblock}
\end{textblock}
\end{frame}


%%%%%%%%%%%%%%%%%%%%%%%%%%%%%%%%%%%%%%%%%%%%%%%%%%%%%%%%%%%%%%%%%%
%%%%%%%%%%%%%%%%%%%%%%%%%%%%%%%%%%%%%%%%%%%%%%%%%%%%%%%%%%%%%%%%%%
\section{Dielectrics}
\begin{frame}{Uniform Plane Waves: Dielectrics}
\begin{textblock}{7.5}(0.5,3.5) 
For dielectrics $\lp \sigma=0, \varepsilon=\varepsilon_r \varepsilon_0, \mu=\mu_0 \rp$:
\begin{small}
\visible<0|handout:1>{
\alert{
	\eq{
		\gamma &= \sqrt{j\omega\mu_0 \lp 0 + j\omega\varepsilon_r\varepsilon_0\rp} \\
		&= \sqrt{-\omega^2 \mu_0 \varepsilon_r\varepsilon_0} \\
		&= j\omega\sqrt{\mu_0\varepsilon_r\varepsilon_0} 			
			= \dfrac{j\omega\sqrt{\varepsilon_r}}{c} \\
			&= \alpha + j\beta \\
		\eta &= \dfrac{j\omega\mu_0}{0 + j\omega\varepsilon_r\varepsilon_0} = \sqrt{\dfrac{\mu_0}{\varepsilon_r\varepsilon_0}} \\
		&= \dfrac{120\pi}{\sqrt{\varepsilon_r}} \\
		u_p &= \dfrac{\omega}{\beta} = \dfrac{\omega}{\omega\sqrt{\mu_0 \varepsilon_0\varepsilon_r}}  = \dfrac{c}{\sqrt{\varepsilon_r}}
	} % end eq
	} % end alert
} % end visible
	
\end{small}
\end{textblock}


\begin{textblock}{7.5}(8.5,2)
\begin{myblock}[0.95]{Free Space}
\visible<0|handout:1>{
\vspace{-12pt}
\alert{
\eq{
	\eta &= \dfrac{120\pi}{\sqrt{1}} = 120\pi = 377\ohms \\
	u_p &= \dfrac{c}{\sqrt{\varepsilon_r}} = \dfrac{c}{\sqrt{1}} = c
} % end eq
} % end alert
\vspace{-12pt}
} % end visible
\end{myblock}

\vspace{-18pt}

\begin{myblock}[0.95]{Silicon, $\varepsilon_r = 11.6$}
\visible<0|handout:1>{
\vspace{-12pt}
\alert{
\eq{
	\eta &= \dfrac{120\pi}{\sqrt{11.6}} = 110.7\ohms\\
	u_p &=  \dfrac{c}{\sqrt{11.6}} =0.3c
} % end eq
} % end alert
\vspace{-12pt}
} % end visible 
\end{myblock}

\end{textblock}
\end{frame}

%%%%%%%%%%%%%%%%%%%%%%%%%%%%%%%%%%%%%%%%%%%%%%%%%%%%%%%%%%%%%%%%%%
%%%%%%%%%%%%%%%%%%%%%%%%%%%%%%%%%%%%%%%%%%%%%%%%%%%%%%%%%%%%%%%%%%
\section{Conductors}
\begin{frame}{Uniform Plane Waves: Conductors}
\begin{textblock}{7.5}(0.5,3.5) 
For conductors:
\begin{small}
\visible<0|handout:1>{
\alert{
	\eq{
	&\lp \sigma >> \omega\varepsilon \rp \text{ or } \lp \dfrac{\sigma}{\omega\varepsilon} >>1 \rp\\[8pt]
	\gamma &= \sqrt{j\omega\mu\lp \sigma + j\omega\varepsilon_r\varepsilon_0\rp} \approx \sqrt{j\omega \mu\sigma} \\[8pt]
	&= e^{j\pi/4}\sqrt{\omega\mu\sigma}\\[8pt]
	&=\lp\cosp{\dfrac{\pi}{4}} + j\sinp{\dfrac{\pi}{4}}\rp\sqrt{\omega\mu\sigma}  \\[8pt]
	&\implies \alpha = \beta = \sqrt{\dfrac{\omega\mu\sigma}{2}}
}% end eq
} % end alert
} % end visible
\end{small}
\end{textblock}

%% blue vertical separator
%\begin{textblock}{0.5}(8.75,3)
%	{\color{Blue4}{\rule{1pt}{2.8in}}}
%\end{textblock}

\begin{textblock}{7.5}(8.5,3.5)
\begin{small}
\visible<0|handout:1>{
\alert{
\eq{
\eta &= \dfrac{j\omega\mu}{\sigma + j\omega\varepsilon} = \sqrt{\dfrac{j\omega\mu}{\sigma}} \\[8pt]
	& = \sqrt{j}\sqrt{\dfrac{\omega\mu}{\sigma}} = \lp e^{j\pi/2} \rp^{1/2} \sqrt{\dfrac{\omega\mu}{\sigma}} \\[8pt]
	&= e^{j\pi/4} \sqrt{\dfrac{\omega\mu}{\sigma}}
} % end eq
} % end alert
} % end visible
\end{small}
\end{textblock}


\begin{textblock}{14.5}(0.5,12)
\begin{myblock}[0.95]{}
\begin{center}
%\vspace{12pt}
	What does this mean?!?!
%\vspace{12pt}
\end{center}
\end{myblock}
\end{textblock}

\end{frame}
%%%%%%%%%%%%%%%%%%%%%%%%%%%%%%%%%%%%%%%%%%%%%%%%%%%%%%%%%%%%%%%%%%
%%%%%%%%%%%%%%%%%%%%%%%%%%%%%%%%%%%%%%%%%%%%%%%%%%%%%%%%%%%%%%%%%%
\begin{frame}{Uniform Plane Waves: Conductors}
\begin{textblock}{7.5}(0.5,2) 
\begin{myblock}[0.95]{$\phasor{E}\overset{\text{?}}{\leftrightarrow}\phasor{H}$}
%How does this affect the relationship between the electric and magnetic field?
\begin{small}
\visible<0|handout:1>{
\alert{
\vspace{-8pt}
\eq{
	\eta & = \dfrac{E_0^{+}}{H_0^{+}} 
		\implies E_0^{+} = \eta H_0^{+} \\
		\implies& E_0^{+} =  H_0^{+} \underbrace{e^{j\pi/4}}_{\text{phase shift}} \sqrt{\dfrac{\omega\mu}{\sigma}}
} % end eq
} % end alert
} % end visible
\end{small}
\end{myblock}
\end{textblock}

\begin{textblock}{7.5}(8.25,2)
\begin{myblock}[0.95]{Dispersion}
	\begin{small}
%		What does $u_p$ look like?
		\visible<0|handout:1>{
		\vspace{12pt}
		\alert{
		\eq{
			u_p &=\dfrac{\omega}{\beta} = \dfrac{\omega}{\sqrt{\frac{\omega\mu\sigma}{2}}} = \sqrt{\dfrac{2\omega}{\mu\sigma}} 
		} % end eq
		} % end alert
		\vspace{12pt}
		} % end visible 
	\end{small}
\end{myblock}
\end{textblock}

\begin{textblock}{14.5}(0.75,9.5)
\begin{myblock}[0.99]{Skin Depth $\sim e^{-\alpha z}$}
	\begin{small}
		Define $\delta$ as skin depth, where $\phasor{E} (z=\delta) = E_0 e^{-1} $:
		\visible<0|handout:1>{
		\vspace{-12pt}
		\alert{
		\eq{
			&\delta = \dfrac{1}{\alpha} = \handoutgap{1in}{\sqrt{\dfrac{2}{\omega\mu\sigma}}}
		} % end eq
		} % end alert
		\vspace{-12pt}
		} % end visible 
	\end{small}
\end{myblock}
\end{textblock}

\end{frame}